\section{Proofs}
\label{app:proofs}

\subsection{Proof of Lemma~\ref{lem:binomial}}

Let $K\sim\mathrm{Binomial}(B,q)$. Two elementary upper bounds suffice. First,
\[
h_B(q)\le\Prob(K>0)\le Bq,
\]
so $qh_B(q)\le Bq^2$. Second, $1/k\le2/(k+1)$ for $k\ge1$, and
\begin{align}
\E\frac1{K+1}
&=\sum_{k=0}^B\frac1{k+1}\binom{B}{k}q^k(1-q)^{B-k}\\
&=\frac{1-(1-q)^{B+1}}{(B+1)q}
\le\frac1{(B+1)q}.
\end{align}
Hence $qh_B(q)\le2/(B+1)$.

For the lower bound, first suppose $Bq\le1$. Then
\begin{align}
h_B(q)
&\ge\Prob(K=1)
=Bq(1-q)^{B-1}\\
&\ge Bq(1-1/B)^{B-1}
\ge e^{-1}Bq,
\end{align}
with the $B=1$ case interpreted directly. Therefore
$qh_B(q)\ge e^{-1}Bq^2$.

Now suppose $Bq\ge1$. Markov's inequality gives
$\Prob(K>2Bq)\le1/2$, while
$\Prob(K=0)=(1-q)^B\le e^{-Bq}\le e^{-1}$. Therefore
\[
\Prob(1\le K\le2Bq)\ge\frac12-e^{-1}.
\]
On this event, $1/K\ge1/(2Bq)$, so
\[
qh_B(q)\ge\frac{1/2-e^{-1}}{2B}.
\]
Combining the two regimes proves the lemma.

